\section*{Chapter \ref{ch:s2:dispersion-and-correlation} --- Dispersion and Correlation}

\begin{solution}{exo:s2:dispersion-and-correlation:1}
For many equal-weighted members the own-variance term vanishes and $\bar\rho \approx \sigma_I^2/\sigma^2 = 0.04/0.09 = 0.444$.
\end{solution}

\begin{solution}{exo:s2:dispersion-and-correlation:2}
By $\sigma_I/\sum_i w_i\sigma_i = 0.2/0.3 = 0.667$.
\end{solution}

\begin{solution}{exo:s2:dispersion-and-correlation:3}
On index weights the book is long more member variance than it is short index variance, so it gains when volatility rises broadly (crashes) and pays
the members' premium in calm months: small regular losses, rare gains.
\end{solution}

\begin{solution}{exo:s2:dispersion-and-correlation:4}
$0.12 - 0.83 = -0.71$ per unit. The strike reflected a calm month's expectations: implied correlation is low when markets are calm, which is when crashes
that raise correlation begin.
\end{solution}

\begin{solution}{exo:s2:dispersion-and-correlation:5}
The premium is measured on the index and all members; trading it needs options on every member, whose spreads and hedging costs are large relative to
the premium: a limit to arbitrage, which is also why the premium can persist.
\end{solution}

\begin{solution}{exo:s2:dispersion-and-correlation:6}
Index options are cheap relative to single-stock options: members are expected to move apart. Selling correlation from there has little premium left and
a large loss if members start moving together, as they do in a sell-off.
\end{solution}

\begin{solution}{exo:s2:dispersion-and-correlation:7}
With the members' premium equal to the index's, implied correlation averages 0.258 against a realised 0.252: the correlation premium nearly vanishes
(0.006). The correlation swap's Sharpe ratio falls to 0.17, the vega-neutral book loses ($-1.84$), and the book on index weights loses every month
($-17.3$), because it pays the full variance premium on more member variance than it sells on the index.
\end{solution}

\begin{solution}{exo:s2:dispersion-and-correlation:8}
Winning in 77\% of months is the shape of a short-volatility trade: the losses come in the others, and the first crash cost the correlation swap six
monthly standard deviations. Leverage chosen on the frequency of wins would not survive the size of the losses; size it by a crash-month stress.
\end{solution}

\begin{solution}{pb:s2:dispersion-and-correlation:1}
\begin{enumerate}
\item $\bar\rho = (\sigma_I^2 - \sum w_i^2\sigma_i^2)/\sum_{i \ne j} w_i w_j\sigma_i\sigma_j$: the single correlation that reconciles index and member
volatilities.
\item The same formula on realised variances; implied minus \omterm{def:s2:dispersion-and-correlation:realised}{realised correlation}.
\item From SPX implied volatility and the implied volatilities of the 50 largest members, with Markowitz's decomposition.
\item Correlation risk is priced and explains option returns; a strategy selling it has a high alpha without frictions and none with them.
\item 37.0 on average, 10.8 to 69.7 for the middle 90\%, 96.6 at the peak (24 October 2008), 70.4 and 72.5 in October 2008 and March 2020.
\item Cboe's S\&P 500 dispersion index: 25.5 on average since 2014, 58.9 at the peak on 18 March 2020.
\item Little premium is left and the loss if correlation rises is large.
\item A smoothed \omterm{def:s2:dispersion-and-correlation:realised}{realised correlation} is in most cases slightly below the implied.
\item Members' notionals on index weights, or scaled to match the index's vega.
\item $-0.02$, 0.65 and 1.74; skewness 2.67, $-5.87$ and $-1.67$.
\item It pays the members' premium on more variance than it sells.
\item It pays implied minus \omterm{def:s2:dispersion-and-correlation:realised}{realised correlation} directly, with no variance-level exposure.
\item \textbf{Named result.} The vega-neutral dispersion book earns the correlation premium (implied 0.311 against realised 0.252) at a Sharpe ratio of
0.65, and loses 11.3 monthly standard deviations when correlation jumps from 0.12 to 0.83 in the first crash; the correlation swap earns 1.74 and loses
0.71 of a unit.
\item 0.12 to 0.83, 0.12 to 0.57, 0.73 to 0.84.
\item Implied correlation was already high when it began.
\item By the loss in a crash month, with correlation going to one.
\item Single-stock option costs, crash months, \omterm{def:s2:dispersion-and-correlation:realised}{realised correlation} as the contract defines it.
\item Vega-weighted dispersion.
\item Chapter~1 sold volatility; dispersion sells the part of index volatility that comes from correlation.
\item Insurance against members moving together.
\end{enumerate}
\end{solution}

\begin{solution}{iq:s2:dispersion-and-correlation:1}
The single average correlation that makes the index's implied variance equal to the weighted sum of members' implied variances and covariances:
$(\sigma_I^2 - \sum w_i^2\sigma_i^2)/\sum_{i \ne j} w_i w_j\sigma_i\sigma_j$ with implied volatilities.
\end{solution}

\begin{solution}{iq:s2:dispersion-and-correlation:2}
Sell index options or variance and buy members' options or variance, weighted by vega, theta or index weights; it earns the correlation premium and
loses when members move together; it also carries single-stock event risk, member liquidity and hedging costs.
\end{solution}

\begin{solution}{iq:s2:dispersion-and-correlation:3}
Index options insure against market-wide falls, when correlation rises and diversification fails; investors pay for that insurance, and index option
supply (overwriting) is smaller than demand, while single-stock options are often supplied by overwriters.
\end{solution}

\begin{solution}{iq:s2:dispersion-and-correlation:4}
Apply historical crash months (2008, 2020) and a hypothetical move of implied and \omterm{def:s2:dispersion-and-correlation:realised}{realised correlation} to one, with members' volatilities up; check the
loss against limits and the liquidity of unwinding the single-stock legs.
\end{solution}

\begin{solution}{iq:s2:dispersion-and-correlation:5}
Index and member option surfaces, index weights and membership changes, corporate events (earnings, mergers) by member, and borrow for hedges; stale
single-stock quotes and membership errors distort implied correlation.
\end{solution}

\begin{solution}{iq:s2:dispersion-and-correlation:6}
With $w_i = 1/N$ and equal $\sigma$, $\sigma_I^2 = \sum_i\sigma^2/N^2 + \sum_{i \ne j}\rho\sigma^2/N^2 = \sigma^2/N + \rho\sigma^2(N - 1)/N = \sigma^2(\rho +
(1 - \rho)/N)$. For large $N$, $\rho \approx \sigma_I^2/\sigma^2$.
\end{solution}
